\makeglossaries

\newacronym{JIT}{JIT}{Just--In--Time compiler}
\newacronym{DSL}{DSL}{Domain--Specific Language}
\newacronym{AST}{AST}{Abstract Syntax Tree}
\newacronym{IR}{IR}{Interal Representation}
\newacronym{GCC}{GCC}{GNU Compiler Colletion}

\newglossaryentry{dsl-gl}{
    name=domain--specific language,
    plural=domain--specific languages,
    description={Szakterület--specifikus nyelv. Olyan programozási nyelv, amelyik egy konkrét probléma megoldására van specializálva, akár annyira, hogy nem is képes általános programozási problémákat megoldani.}
}
\newglossaryentry{sexpr}{
    name=S--expr,
    first=symbolic expression ({\it S--expression}),
    description={Symbolic expression. Olyan egymásba ágyazott listákat leíró szintaktia, amelyet többek között a LISP nyelvcsalád használ, mind kód, mind adat tárolására}
}
\newglossaryentry{prefix}{
    name=prefix kifejezés,
    description={Olyan kifejezés, amelyikben az elvégzendő művelet a kifejezés legelején található. \pl egy LISP \gls{sexpr}, {\tt (+ 1 2)}},
    see={infix,sexpr}
}
\newglossaryentry{infix}{
    name=infix kifejezés,
    description={Olyan kifejezés, amelyikben az elvégzendő művelet a kifejezés közepén található. \pl a legtöbb szokásos matematikai művelet jelölés, $1 + 2$},
    see={prefix}
}
\newglossaryentry{pure}{
    name=tiszta,
    first=tiszta ({\it pure}),
    description={Olyan függvény, amelyik megfelel a matematikai függvény definíciójának: adott bemenetekre mindig ugyan azt az eredményt adja és nincs semmiféle mellékhatása. Tehát olyan megfigyelhető változást nem eredményez, amelyik nem a kimeneti értéke}
}
\newglossaryentry{precedence}{
    name=precedencia,
    first=precedencia ({\it precedence}),
    description={Két operátor egy kifejezésben használatakor, azok kiértékelési sorrendjét megadó sorrendezése minden operátornak}
}
\newglossaryentry{associativity}{
    name=asszociativitás,
    first=asszociativitás ({\it associativity}),
    description={Két azonos operátor egymásutánjakor a zárójelezés irányát meghatározó tulajdonsága egy operátornak. Lehet jobb vagy bal}
}
\newglossaryentry{token}{
    name=token,
    plural=tokenek,
    description={Elemi egységei egy--egy nyelvtannak, amelyek a szöveg egy folyamatos részéhez rendelnek szemantikus értelmet. Tartalmazhat valamilyen formában értelmezett részeket is, \pl egy szám értéke már nem csak szöveg, hanem szám formában is.}
}
\newglossaryentry{memory-mapping}{
    name=memóriába--lapozás,
    first=memóriába--lapozás (memory mapping),
    description={Egy folyamat, amelyen keresztül az operációs rendszert felhasználva egy fájlt tartalmát memóriában elérhetővé alakítjuk.}
}
\newglossaryentry{eager}{
    name=szorgos,
    first=szorgos ({\it eager}),
    description={Olyan kiértékelés, amelyikben az összes kifejezés azonnal kiértékelésre kerül, attól függetlenül, hogy használja-e annak a végeredményét másik kifejezés}
}
\newglossaryentry{lazy}{
    name=lusta,
    first=lusta ({\it lazy}),
    description={Olyan kiértékelés, amelyikben egy kifejezés csak akkor értékelődik ki ténylegesen, ha van olyan másik kifejezés, amelyik kiértékelődik, és használja az első kifejezés eredményét.}
}
\newglossaryentry{closure}{
    name=closure,
    description={Olyan blokk, amelyik saját számításaihoz egy rajta kívül definiált, de nem globláis függvényt használ.}
}